\documentclass[10pt,a4paper]{amsart}
\usepackage[margin=19mm]{geometry}
\usepackage{amsmath,amssymb,mathtools}
\usepackage[scr=rsfs,cal=euler]{mathalfa}
\usepackage{xfrac}
\usepackage[unicode,hidelinks,bookmarks=false]{hyperref}
\newcommand{\ie}{i.e.\ }
\newcommand{\cf}{cf.\ }
\newcommand{\resp}{respectively\ }
\newcommand{\Subsection}{\subsection}
\providecommand{\nobreakdash}{\nobreak\hbox{-}}
\setlength{\emergencystretch}{2em}
\multlinegap0pt
\usepackage{xfrac}
\usepackage{amssymb}
\usepackage{xcolor}
\usepackage[shortlabels]{enumitem}
\usepackage{tikz}
\usepackage{cancel}
\newtheorem{proposition}{Proposition}
\newcommand{\R}{\mathbb{R}}
\newcommand{\rank}{\operatorname{rank}}
\title{Sobolev mappings of Euclidean space and product structure}
\author{Bruce Kleiner, Stefan M\"uller, L\'{a}szl\'{o} Sz{\'{e}}kelyhidi, Jr., Xiangdong Xie}
\date{Source: arXiv:2603.05979v1 (CC BY 4.0)}
\begin{document}
\maketitle

\noindent\emph{Source and licence.} Sobolev mappings of Euclidean space and product structure. Version arXiv:2603.05979v1 (CC BY 4.0), distributed by arXiv under the Creative Commons Attribution 4.0 International licence, http://creativecommons.org/licenses/by/4.0/, as recorded in the arXiv licence metadata for this exact version and shown on the abstract page https://arxiv.org/abs/2603.05979v1. Full licence notice: \texttt{licenses/arxiv-2603.05979-CC-BY-4.0.txt}.\par\smallskip

\noindent\textit{Editorial scope.} Counted unit: the article's proposition pr:T5\_permutation (source0.tex lines 1688--1722). The article's complete original proof of it is delivered with it (source0.tex lines 1823--2077, one \texttt{proof} environment of 255 lines, which the article places away from the statement). No mathematics is written, repaired or re-derived: the statement and the proof are the article's own lines, in the article's own notation, and no retained line is substituted or re-flowed.  Retained uncounted as the source-local context the proof needs: lines 1641--1643; lines 1768--1779; lines 1770--1772; lines 1811--1818.  Nothing was omitted from any retained span, so the package carries no omission record.  No retained line contains a citation, so the wrapper carries no bibliography and no bibliography entry.  No retained line contains a figure environment, an image-inclusion command or drawing code, so no image is delivered and the package declares no asset dependency.  Editorial formatting only, in addition: an AMS \texttt{amsart} preamble that restates the article's own declarations and macros used by the retained lines, namely the environments \texttt{proposition} on separate counters, together with the standard packages those lines need.  The retained environments print in this document as Proposition 1 in order of appearance, whereas the article prints the same items with the numbering of its own full document.  The complete native source of the article as distributed in the arXiv e-print for this exact version is bundled unchanged as \texttt{sources/195828/source0.tex}.
\begin{equation}\label{T5e:mu*}
\mu_*=1+\frac{\beta}{2}+\frac{1}{2}\sqrt{\beta^2+4\beta}.
\end{equation}

\begin{proposition} \label{pr:criterion_rank3}  Assume that $\{X_1,\dots,X_5\}\subset\R^{2\times 2}$ is affine non-degenerate, meaning that the affine subspace of $\R^{2\times 2}$ spanned by these $5$ matrices is $4$-dimensional. Then the large $T_5$ property is 
equivalent to the condition that there exist $\sigma_1,\sigma_2,\sigma_3$ $T_5$-configurations and furthermore 
\begin{equation}\label{T5e:rank}
	\textrm{rank }B^{(k)}=3\textrm{ for all }k=1,\dots,5,
\end{equation}
where
$$
B^{(k)}_{ij}=\begin{cases} \lambda_j^{\sigma_i}& \sigma_i^{-1}(k)<\sigma_i^{-1}(j),\\
0&k=j,\\
\mu^{\sigma_i}\lambda_j^{\sigma_i}&	\sigma_i^{-1}(k)>\sigma_i^{-1}(j).\end{cases}
$$
\end{proposition}

\begin{equation}
	\textrm{rank }B^{(k)}=3\textrm{ for all }k=1,\dots,5,
\end{equation}

\begin{proposition}\label{T5p:ex}
Let  $a=(c-\tfrac1c)^2>0$,  $b=c+\tfrac1c-2>0$ and
\begin{equation*}
\sigma_1=[1\, 2\, 3\, 5\, 4],\quad 	\sigma_2=[1\, 2\, 4\, 5\, 3],\quad 	\sigma_3=[1\, 2\, 5\, 3\, 4].
\end{equation*}
If $ab>8$, then the set $\{X_1,\dots,X_5\}$ is a large $T_5$ set. More precisely, in this case the permutations $\sigma_j$, $j=1,2,3$ correspond
to $T_5$-configurations and the associated rank-one directions $\{P_i^{\sigma_j}-X_i:\,j=1,2,3\}$ are linearly independent for each $i=1,\dots,5$.
\end{proposition}

\begin{proposition}   \label{pr:T5_permutation} 
(i) The tuple
$X^\sigma := \left(X_{\sigma(1)},X_{\sigma(2)},X_{\sigma(3)},X_{\sigma(4)},X_{\sigma(5)}\right)$ is a $T_5$ configuration if and only if 
there exist $\mu^{\sigma}>1$ and $\lambda_1^{\sigma},\dots,\lambda_5^{\sigma}>0$ such that
$$
A^{\sigma,\mu^{\sigma}}\lambda^{\sigma}=0,
$$
where
$$
A_{ij}^{\sigma,\mu}=\begin{cases}\det(X_i-X_j)&\sigma^{-1}(i)<\sigma^{-1}(j),\\ 
0&i=j,\\
\mu\det(X_i-X_j)&\sigma^{-1}(i)>\sigma^{-1}(j).\end{cases} 
$$
Here  $\sigma^{-1}$ denotes  the inverse permutation of $\sigma$. 

(ii) If $X^\sigma$ is a $T_5$ configuration then the set of inner points is given by
$\{ P^\sigma_1, \ldots P^\sigma_5\}$ where   
\begin{equation} \label{T5e:Pki}
P_k^\sigma = \sum_{i=1}^5  \xi^{(\sigma, k)}_i X_i
\end{equation}
with  $\xi^{(\sigma, k)} = (\sum_{i} v_i^{(\sigma,k)})^{-1} v^{(\sigma,k)}$, 
\begin{equation}\label{T5e:vki} 
v^{(\sigma,k)}_i = \begin{cases}\lambda^\sigma_i & \sigma^{-1}(i)\geq \sigma^{-1}(k)\\ \mu^\sigma\lambda^\sigma_i & \sigma^{-1}(i)< \sigma^{-1}(k), \end{cases}
\end{equation}
and $v^{(\sigma, k)} = \sum_{i=1}^5 v^{(\sigma, k)}_i$.
In particular $X_k$ is rank-one connected to $P^\sigma_k$. 

(iii) Set $A = A^{\sigma,1}$ and 
\begin{equation}\label{e:permuteddeterminants}
\alpha^\sigma := A_{\sigma(1) \sigma(2)}  A_{\sigma(2) \sigma(3)} A_{\sigma(3) \sigma(4)} A_{\sigma(4) \sigma(5)}  A_{\sigma(5) \sigma(1)}.
\end{equation}
Then $ \mu \mapsto \det A^{\sigma, \mu}$  has a zero  $\mu^* > 1$ if and only if
$$ \beta^\sigma := -  \frac{\det A}{\alpha^\sigma} > 0.$$
If this condition holds then $\mu^*$ is given by  \eqref{T5e:mu*} with $\beta$ replaced by $\beta^\sigma$.
\end{proposition}

\begin{proof}[Proof of Proposition \ref{T5p:ex}]
Let $X^\sigma = (X_{\sigma(1)}, \ldots, X_{\sigma(5)})$. 
We first show that $X^{\sigma_1}$, $X^{\sigma_2}$ and $X^{\sigma_3}$ are
$T_5$ configurations by using the criterion in Proposition~\ref{pr:T5_permutation}.



To compute $A = A^{\sigma, 1}$ we observe that 
 $\det X_i=1$ for all $i$. 
  Thus
\begin{equation*}
A=\begin{pmatrix}
	0&-a&2&2&-b\\ -a&0&2&2&-b\\2&2&0&-a&2\\2&2&-a&0&2\\-b&-b&2&2&0
\end{pmatrix}	
\end{equation*}
with $a=(c-\tfrac1c)^2>0$, and $b=c+\tfrac1c-2>0$. 
A direct calculation gives
$$
\det A=2a^2(ab^2+4a-16b).
$$
Since $a=b(c^{\sfrac12}+c^{-\sfrac12})^2> 4b$, we see that $\det A>0$. 
Moreover, using the notation from \eqref{e:permuteddeterminants},
\begin{eqnarray*} \alpha^{\sigma_1} &=& A_{12} A_{23} A_{35} A_{54} A_{41} = - 16 a, \\
 \alpha^{\sigma_2} &=& A_{12} A_{24} A_{45} A_{53} A_{31} = - 16 a, \\
  \alpha^{\sigma_3} &=& A_{12} A_{25} A_{53} A_{34} A_{41} = - 4 a^2 b.
  \end{eqnarray*}
Hence $\beta^{\sigma_i}>0$ and therefore the $\mu^{(i)}$ defined by formula \eqref{T5e:mu*}  with $\beta$ replaced by
$\beta^{\sigma_i}$ satisfy $\mu^{(i)}>1$. 
Note also that, since $ab>8$, $\beta^{\sigma_1}=\beta^{\sigma_2}>\beta^{\sigma_3}$ 
and consequently $\mu^{(1)}=\mu^{(2)}>\mu^{(3)}$. In  the following, let us denote
$$
\eta=\mu^{(1)}=\mu^{(2)},\quad \nu=\mu^{(3)}.
$$

It remains to check that the kernels of the matrices $A^{\sigma_i,\mu^{(i)}}$ intersect the 1st octant $(0,\infty)^5$.
 Let us first consider the permutation $\sigma_1=[1 2 3 5 4]$. Then 
$$
A^{\sigma_1,\eta}=\begin{pmatrix}
	0&-a&2&2&-b\\ -a\eta &0&2&2&-b \\ 2\eta &2\eta &0&-a&2\\ 2\eta &2\eta &-a\eta &0&2\eta\\-b\eta &-b\eta &2\eta &2 &0
\end{pmatrix}.	
$$ 
Subtracting the first row from the second and subsequently subtracting appropriate multiples of the second row from the others, we obtain
$$
\tilde A^{\sigma_1,\eta}=\begin{pmatrix}
	0&-a&2&2&-b\\ 
	\eta &-1&0&0&0\\ 
	2\eta(1+\eta) &0&0&-a&2\\ 
	2\eta(1+\eta) &0&-a\eta &0&2\eta\\
	-b\eta(1+\eta) &0 &2\eta &2 &0
\end{pmatrix},
$$ 
and a further row reduction results in the matrix
$$
  \tilde{ \tilde A}^{\sigma_1, \eta}  = \begin{pmatrix}
	0&-a&2&2&-b\\ 
	\eta &-1&0&0&0\\ 
	\lambda^{(1)}_3 &0&-1&0&0\\ 
	\lambda^{(1)}_4 &0&0&-1&0\\
	\lambda^{(1)}_5 &0 &0&0&-1
\end{pmatrix},
$$ 
where $\lambda^{(1)}\in \R^5$ is given by
$$
\lambda^{(1)}=\begin{pmatrix}1&\eta &\tfrac{2}{a}\left((\tfrac{ab}{4}-1)\eta+1\right)&\tfrac{2\eta}{a}\left(\tfrac{ab}{4}-1+\eta\right)&\left(\frac{ab}{4}-2\right)\eta\end{pmatrix}^T
$$ 
In particular, $( \tilde {\tilde A}^{\sigma_1,\eta}\lambda^{(1)})_i=0$ for $i=2,\dots,5$.

Clearly the last four rows of $\tilde {\tilde A}^{\sigma_1,\eta}$
are linearly independent. Since $\det \tilde {\tilde A}^{\sigma_1,\eta} = 0$
it follows that the first row of $\tilde {\tilde A}^{\sigma_1,\eta}$ is a linear combination
of the last four rows. Thus $\tilde {\tilde A}^{\sigma_1,\eta} \lambda^{(1)} = 0$
and hence $A^{\sigma_1, \eta} \lambda^{(1)} = 0$. Since $\rank A^{\sigma_1, \eta} = 4$ it follows that $\lambda^{(1)}$ is the vector which generates the kernel of 
$ A^{\sigma_1, \eta} $. 
Since $a,b>0$, $\eta>1$ and $ab>8$, we see $\lambda^{(1)}_i>0$ for all $i$, so that $\sigma_1$ indeed corresponds to a $T_5$-configuration.


\smallskip


Concerning the case $\sigma_2=[1 2 4 5 3]$ we note that
$$A^{\sigma_2,\eta}=\begin{pmatrix}
	0&-a&2&2&-b\\ -a\eta &0&2&2&-b \\ 2\eta &2\eta &0&-a \eta &2 \eta \\ 2\eta &2\eta &-a &0&2\\-b\eta &-b\eta &2 &2 \eta &0
\end{pmatrix}.
$$
Comparing this  expression with 
$$ A^{\sigma_1,\eta}=\begin{pmatrix}
	0&-a&2&2&-b\\ -a\eta &0&2&2&-b \\ 2\eta &2\eta &0&-a&2\\ 2\eta &2\eta &-a\eta &0&2\eta\\-b\eta &-b\eta &2\eta &2 &0
\end{pmatrix}.	
$$
we see that $A^{\sigma_2, \eta}$ is obtained from $A^{\sigma_1, \eta}$ for swapping 
the 3rd and 4th row and the 3rd and 4th column.



  Hence
$$
\lambda^{(2)}=\begin{pmatrix}1&\eta &\tfrac{2\eta}{a}\left(\tfrac{ab}{4}-1+\eta\right)&\tfrac{2}{a}\left((\tfrac{ab}{4}-1)\eta+1\right)&\left(\frac{ab}{4}-2\right)\eta\end{pmatrix}^T
$$
is the vector generating the $1$-dimensional kernel of $A^{\sigma_2,\eta}$, and, as above, we see that $\lambda^{(2)}_i>0$ for all $i$ under the conditions of the proposition.  

\smallskip

Finally, let us look at $\sigma_3=[1 2 5 3 4]$. Here
$$
A^{\sigma_3,\nu}=\begin{pmatrix}
	0&-a&2&2&-b\\ -a\nu &0&2&2&-b \\ 2\nu &2\nu &0&-a&2\nu\\ 2\nu &2\nu &-a\nu &0&2\nu\\-b\nu &-b\nu &2 &2 &0
\end{pmatrix}.	
$$ 
Proceeding with row-reduction as above, we obtain first
$$
\tilde A^{\sigma_3,\nu}=\begin{pmatrix}
	0&-a&2&2&-b\\ 
	\nu &-1&0&0&0\\ 
	2\nu(1+\nu) &0&0&-a&2\nu\\ 
	2\nu(1+\nu) &0&-a\nu &0&2\nu\\
	-b\nu(1+\nu) &0 &2 &2 &0
\end{pmatrix},
$$ 
and a further row reduction results in the matrix
$$
  \tilde{\tilde A}^{\sigma_3, \nu} 
=
\begin{pmatrix}   	0&-a&2&2&-b\\ 
	\nu &-1&0&0&0\\ 
	\lambda^{(3)}_3 &0&-1&0&0\\ 
	\lambda^{(3)}_4 &0&0&-1&0\\
	\lambda^{(3)}_5 &0 &0&0&-1
\end{pmatrix},
$$ 
where
$$
\lambda^{(3)}=\begin{pmatrix}1&\nu &\tfrac{b}{2}\nu &\tfrac{b}{2}\nu^2&\left(\frac{ab}{4}-1\right)\nu-1\end{pmatrix}^T.
$$

Arguing as before, we deduce that $A^{\sigma_3,\nu}\lambda^{(3)}=0$, and furthermore $\lambda^{(3)}_i>0$ for all $i$ since $a,b>0$, $\nu>1$ and $ab>8$. Therefore $\sigma_3$ also corresponds to a $T_5$ configuration.

\bigskip






In view of Proposition~\ref{pr:criterion_rank3} it only
  remains to check the rank condition in \eqref{T5e:rank}. To this end it suffices to show that for each of the $3\times 5$ matrices $B^{(k)}$ there exists a non-vanishing $3\times 3$ subdeterminant. A judicious choice of the relevant columns in each $B^{(k)}$ leads one to look at 
$$
B^{(1)}_{[2 4 5]},\,B^{(2)}_{[1 4 5]},\,
 B^{(3)}_{[124]},\,
B^{(4)}_{[1 2 3]},\,B^{(5)}_{[1 2 3]},
$$
where $B^{(k)}_{[l m n]}$ denotes the $3\times 3$ matrix formed by restricting $B^{(k)}$ to columns $[l m n ]$. 
Elementary calculations lead to
\begin{align*}
\det B^{(1)}_{[2 4 5]} 
&
 = \det \begin{pmatrix} \lambda^{(1)}_2   &  \lambda^{(1)}_4 & \lambda^{(1)}_5 \\
  \lambda^{(2)}_2 &  \lambda^{(2)}_4 & \lambda^{(2)}_5\\
   \lambda^{(3)}_2   &  \lambda^{(3)}_4 & \lambda^{(3)}_5
\end{pmatrix}
\\
 &= \det \begin{pmatrix}
	\eta &\tfrac{2\eta}{a}\left(\tfrac{ab}{4}-1+\eta\right)&\left(\frac{ab}{4}-2\right)\eta\\
    \eta &\tfrac{2}{a}\left((\tfrac{ab}{4}-1)\eta+1\right)&\left(\frac{ab}{4}-2\right)\eta\\
    \nu &\tfrac{b}{2}\nu^2&\left(\frac{ab}{4}-1\right)\nu-1
\end{pmatrix}\\
&=\det\begin{pmatrix}
	\eta &\tfrac{2\eta}{a}\left(\tfrac{ab}{4}-1+\eta\right)&\left(\frac{ab}{4}-2\right)\eta\\
    0 &\tfrac{2}{a}\left(1-\eta^2\right)&0\\
    \nu &\tfrac{b}{2}\nu^2&\left(\frac{ab}{4}-1\right)\nu-1
\end{pmatrix}\\
&=-\tfrac{2}{a}(\eta^2-1)\eta(\nu-1),\\
  B^{(2)}_{[1 4 5]} 
 & =
 \begin{pmatrix} \eta \lambda^{(1)}_1 &  \lambda^{(1)}_4 & \lambda^{(1)}_5 \\
\eta  \lambda^{(2)}_1 &  \lambda^{(2)}_4 & \lambda^{(2)}_5\\
 \nu   \lambda^{(3)}_1 &   \lambda^{(3)}_4 & \lambda^{(3)}_5
\end{pmatrix}  
=  B^{(1)}_{[245]},  \quad  \text{and thus}  \\
\det   B^{(2)}_{[1 4 5]} & = \det  B^{(1)}_{[245]}.
 \end{align*}
 

 
 Moreover, 
 \begin{align*}
\det B^{(3)}_{[124]} &=
%
\det  \begin{pmatrix} \eta \lambda^{(1)}_1  & \eta   \lambda^{(1)}_2 &   \lambda^{(1)}_4 &  \\
\eta  \lambda^{(2)}_1 & \eta \lambda^{(2)}_2 &   \eta  \lambda^{(2)}_4 &  \\
 \nu   \lambda^{(3)}_1 & \nu \lambda^{(3)}_2 &   \lambda^{(3)}_4 
\end{pmatrix}
  \\
%
&= 
\det\begin{pmatrix}
	\eta &\eta^2&\tfrac{2\eta}{a}\left(\tfrac{ab}{4}-1+\eta\right)\\
    \eta &\eta^2&\tfrac{2\eta}{a}\left((\tfrac{ab}{4}-1)\eta+1\right)\\
    \nu &\nu^2&   \tfrac{b}{2}\nu^2
\end{pmatrix}  \\
%
&= \eta^2 \nu
 \det\begin{pmatrix} 
	1 &\eta&\tfrac{2}{a}\left(\tfrac{ab}{4}-1+\eta\right)\\
    1 &\eta&\tfrac{2}{a}\left((\tfrac{ab}{4}-1)\eta+1\right)\\
    1 &\nu&  \tfrac{b}{2}\nu
\end{pmatrix}
\\
&=
\eta^2 \nu
 \det\begin{pmatrix} 
	1 &\eta&\tfrac{2}{a}\left(\tfrac{ab}{4}-1+\eta\right)\\
    0 & 0 &\tfrac{2}{a}     \left((\tfrac{ab}{4}-1)\eta+1     - \tfrac{ab}{4} + 1 - \eta) \right)\\
    1 &\nu&  \tfrac{b}{2}\nu
\end{pmatrix}
\\
&=-\frac{2}{a}\eta^2\nu(\tfrac{ab}{4}-2)(\eta-1)(\nu-\eta)
\end{align*}

This calculation shows in particular that the determinant of three times three matrix  $B^{(3)}_{[124]}$
does not depend on the $33$ entries of the matrix. Using this fact we obtain in the same way
%
%
\begin{align*}
\det B^{(4)}_{[1 2 3]}
&=
%
\det \begin{pmatrix} \eta \lambda^{(1)}_1  & \eta   \lambda^{(1)}_2 &  \eta  \lambda^{(1)}_3 \\
\eta  \lambda^{(2)}_1 & \eta \lambda^{(2)}_2 &    \lambda^{(2)}_3 \\
 \nu   \lambda^{(3)}_1 & \nu \lambda^{(3)}_2 &   \nu \lambda^{(3)}_3 \end{pmatrix}
\\
%
&=\det\begin{pmatrix}
	\eta &\eta^2&\tfrac{2\eta}{a}\left((\tfrac{ab}{4}-1)\eta+1\right)\\
    \eta &\eta^2&\tfrac{2\eta}{a}\left(\tfrac{ab}{4}-1+\eta\right)\\
    \nu &\nu^2&   \tfrac{b}{2}\nu^2
\end{pmatrix}\\
&=    -\det B^{(3)}_{[124]}  \\%
%
\det B^{(5)}_{[1 2 3]}  &=
%
\det  \begin{pmatrix} \eta \lambda^{(1)}_1  & \eta   \lambda^{(1)}_2 &  \eta  \lambda^{(1)}_3  \\
\eta  \lambda^{(2)}_1 & \eta \lambda^{(2)}_2 &    \lambda^{(2)}_3 \\
 \nu   \lambda^{(3)}_1 & \nu \lambda^{(3)}_2 &    \lambda^{(3)}_3 
\end{pmatrix}
\\
%
&=\det\begin{pmatrix}
	\eta &\eta^2&\tfrac{2\eta}{a}\left((\tfrac{ab}{4}-1)\eta+1\right)\\
    \eta &\eta^2&\tfrac{2\eta}{a}\left(\tfrac{ab}{4}-1+\eta\right)\\
    \nu &\nu^2&     \tfrac{b}{2}\nu
\end{pmatrix}\\
&= \det B^{(3)}_{[124]}. \\
\end{align*}
Since we already know that $\eta>\nu$, $\eta,\nu>1$ and $ab>8$ by assumption, we deduce that none of the $5$ determinants above vanishes, thus showing that the rank condition \eqref{T5e:rank} is satisfied. This concludes the proof.
\end{proof}

\end{document}
